\documentclass{beamer}
\usepackage{beamerthemesplit}
\usepackage{wrapfig}
\usetheme{SPbGU}
\usepackage{pdfpages}
\usepackage{amsmath}
\usepackage{cmap} 
\usepackage[T2A]{fontenc} 
\usepackage[utf8]{inputenc}
\usepackage[english,russian]{babel}
\usepackage{indentfirst}
\usepackage{amsmath}
\usepackage{tikz}
\usepackage{multirow}
\usepackage[noend]{algpseudocode}
\usepackage{algorithm}
\usepackage{algorithmicx}
\usetikzlibrary{shapes,arrows}
\usepackage{fancyvrb}
\usepackage{tikz}
\usepackage{pgfplots}
\usetikzlibrary{calc}
\usetikzlibrary{shapes,arrows}
\usetikzlibrary{arrows,automata}
\usetikzlibrary{positioning}
\usepackage[labelformat=empty]{caption}
\usepackage[labelformat=empty]{subcaption}
\beamertemplatenavigationsymbolsempty
\usepackage{multirow}
\usepackage{array}
\newcolumntype{P}[1]{>{\centering\arraybackslash}p{#1}}


\newcommand{\tikzmark}[1]{\tikz[overlay,remember picture] \node (#1) {};}
\def\Put(#1,#2)#3{\leavevmode\makebox(0,0){\put(#1,#2){#3}}}

\newcommand{\ltz}{$< 1$}
\tikzset{
    state/.style={
           rectangle,
           rounded corners,
           draw=black, very thick,
           minimum height=2em,
           inner sep=4pt,
           text centered,
           },
}

\title[СА + НС]{Комбинирование нейронных сетей и синтаксического анализа для предсказания вторичных структур генетических цепочек}
\institute[СПбГУ]{
Санкт-Петербургский государственный университет \\
Кафедра системного программирования }

\author[Лунина Полина]{Лунина Полина Сергеевна, 571 группа \\
  \and  
    {\bfseries Научный руководитель:} доцент, к.ф-м.н. Григорьев С.В.}


\date{10 июня 2020г.}
\begin{document}
{
\begin{frame}[fragile]
  \begin{tabular}{p{2.0cm} p{7.5cm} p{1cm}}
   \begin{center}
      \includegraphics[height=1.5cm]{pics/jetbrainsResearch.pdf}
    \end{center}
    &
    \begin{center}
      \includegraphics[height=1.5cm]{pics/YC_logo.pdf}
    \end{center}
    &
    \begin{center}
      \includegraphics[height=1.5cm]{pics/SPbGU_Logo.png}
    \end{center}
  \end{tabular}
  \titlepage
\end{frame}
}

\begin{frame} \frametitle{Биоинформатика}
\begin{tabular}{cl}  
    \parbox{0.44\linewidth}{
        \begin{itemize}
            \item Задачи
            \begin{itemize}
                \item Распознавание
                \item Классификация
                \item Предсказание вторичных структур
                \item ...
            \end{itemize}
  \onslide<2,3> {\item Формальное задание вторичной структуры}
  \onslide<3>  {\item Вероятностная оценка}
        \end{itemize}
    }
    & \begin{tabular}{l}
        \vspace{-0.8cm}
        \hspace{-0.8cm}
        \includegraphics[width=5.0cm]{pics/16s.pdf}
    \end{tabular}  \\
\end{tabular}
\end{frame}

\begin{frame}{Наш подход}
\begin{itemize}
    \item Задать основные элементы вторичной структуры (стемы) с помощью грамматики 
    \item Для вероятностной оценки использовать нейронные сети
\end{itemize}

\vspace{1.0cm}
\begin{center}
    \includegraphics[width=8.5cm]{pics/molekula.pdf} 
\end{center}
    
\end{frame}

\begin{frame} \frametitle{Наш подход}
\vspace{0.5cm}
\scriptsize
  \begin{tikzpicture}[->,>=stealth']
  \node[state,
        align=left,
        text width = 3.2cm] (grm)
   {
   \textbf{Грамматика}\\
   Описывает стемы вторичной структуры.
   };

  \node[state,
        below of=grm,
        node distance=1.6cm,
        align = left,
        text width = 3.2cm] (sqs)
   {
   \textbf{Цепочки РНК}\\
   Строки над алфавитом $\{A, C, G, U\}$.
   };

  \node[state,
        below right = -0.4cm and 0.4cm of grm,
        align = left,
        text width = 3.0cm](parser)
  {
  \textbf{Парсер}\\
  Извлекает стемы вторичной структуры из входных данных.
  };

  \node[state,
        right of=parser,
        node distance=4.3cm,
        align = left,
        text width = 4.4cm] (mtrx)
  {
\textbf{Матрицы}\\
{\tiny
\[ \left( \begin{array}{cccc}
0 & 1 & 0 & 1\\
0 & 0 & 1 & 0\\
0 & 0 & 0 & 1\\
0 & 0 & 0 & 0
\end{array} \right)\]}
\vspace{-0.2cm}

  Результат работы парсера для входной строки $\omega$ --- булева матрица $M$, где $ M[i,j] = 1 $ $\iff$  \linebreak  $\text{\ttfamily{s}} \xrightarrow{*}{} \omega[i,j]$.
  };


  \node[state,
        below = 1.2cm of mtrx,
        node distance=3cm,
        align = left,
        text width = 4.4cm](vector)
  {
  \textbf{Входные данные для НС}\\
  Изображения, вектора или другой формат, удобный для конкретной задачи.
  };

  \node[state,
        left = 0.5cm of vector,
        align = left,
        text width = 4.4cm](DNN)
  {
  \textbf{Нейронная сеть}\\

  Решает поставленную задачу распознавания, классификации, генерации и т.д.

  };
  
    \node[state,
        left of = DNN,
        node distance=3.7cm,
        align = left,
        text width = 1.4cm](result)
  {
  \textbf{Результат}
  };

  \path (grm.east) edge (parser.-200)
   (sqs.east) edge (parser.-160)
   (parser) edge (mtrx)
   (mtrx) edge (vector)
   (vector) edge (DNN)
   (DNN) edge (result)
   ;

  \end{tikzpicture}

  \tikzmark{zzz}{
  }
  \pause
  \onslide<2>{\tikz[overlay,remember picture]{\draw[draw=red,thick,double,fill opacity=0.2, rounded corners] ($ (zzz) + (0.0,6.4)$) rectangle ($ (zzz) + (3.52,3.56)$);}}

  \onslide<3>{\tikz[overlay,remember picture]{\draw[draw=red,thick,double,fill opacity=0.2, rounded corners] ($ (zzz) + (3.95,5.55)$) rectangle ($ (zzz) + (7.22,4.0)$);}}

  \onslide<4>{\tikz[overlay,remember picture]{\draw[draw=red,thick,double,fill opacity=0.2, rounded corners] ($ (zzz) + (7.55,6.45)$) rectangle ($ (zzz) + (12.22,3.15)$);}}

  \onslide<5>{\tikz[overlay,remember picture]{\draw[draw=red,thick,double,fill opacity=0.2, rounded corners] ($ (zzz) + (7.55,1.86)$) rectangle ($ (zzz) + (12.22,0.35)$);}}

  \onslide<6>{\tikz[overlay,remember picture]{\draw[draw=red,thick,double,fill opacity=0.2, rounded corners] ($ (zzz) + (2.3,1.86)$) rectangle ($ (zzz) + (7.0,0.35)$);}}

  \onslide<7>{\tikz[overlay,remember picture]{\draw[draw=red,thick,double,fill opacity=0.2, rounded corners] ($ (zzz) + (0.1,1.4)$) rectangle ($ (zzz) + (1.8,0.8)$);}}

\end{frame}

\begin{frame}{Идея текущего исследования} 
\begin{itemize}
    \item Парсер находит в цепочке все возможные стемы, однако не все они действительно будут входить в состав вторичной структуры
    \item Хотим сконструировать нейронную сеть, которая отфильтрует лишние контакты между нуклеотидами и предскажет вторичную структуру цепочки
\end{itemize}
\vspace{0.5cm}
\onslide<2>{
\begin{figure}
\centering
\begin{subfigure}{.3\textwidth}
  \centering
  \fbox{\includegraphics[width=.6\linewidth]{pics/structure.png}}
  \caption{Вторичная структура}
\end{subfigure}%
\begin{subfigure}{.3\textwidth}
  \centering
  \fbox{\includegraphics[width=.6\linewidth]{pics/etalon.png}}
  \caption{Матрица контактов}
\end{subfigure}
\begin{subfigure}{.3\textwidth}
  \centering
  \fbox{\includegraphics[width=.6\linewidth]{pics/parsed.png}}
  \caption{Матрица разбора}
\end{subfigure}%
\end{figure}
}
\end{frame}

\begin{frame}\frametitle{Постановка задачи}
\textbf{Цель} --- исследование возможности применения предложенного подхода к задаче предсказания вторичных структур геномных последовательностей

\vspace{5mm\onelineskip}

\textbf{Задачи}
\begin{itemize}
    \item Разработка общей архитектуры решения
    \item Проведение экспериментальных исследований
    \begin{itemize}
        \item Предсказание вторичных структур транспортных РНК с различной длиной цепочки
        \item Исследование возможности предсказания псевдоузлов, невыразимых средствами используемой грамматики
    \end{itemize}
\end{itemize}
\end{frame}

\begin{frame}{Архитектура решения}
\scriptsize
\centering
  \begin{tikzpicture}[->,>=stealth']
  \node[state,
        align=center,
        text width = 2.0cm] (grm)
   {
   \textbf{Грамматика}
   };

  \node[state,
        below of=grm,
        node distance=2.0cm,
        align = center,
        text width = 2.0cm] (sqs)
   {
   \textbf{Цепочки РНК}
   };

  \node[state,
        right = 0.5cm of grm,
        align = center,
        double,
        text width = 1.5cm](parser)
  {
  \textbf{Парсер}
  };
  
  \node[state,
        right = 0.5cm of sqs,
        align = center,
        double,
        text width = 1.5cm](pred)
  {
  \textbf{Генератор в.с.}
  };

  \node[state,
        right = 0.5cm of parser,
        align = center,
        text width = 1.5cm] (mtrx1)
  {
\textbf{Матрицы разбора}\\

};

  \node[state,
        right = 0.5cm of pred,
        align = center,
        text width = 1.5cm](mtrx2)
  {
  \textbf{Матрицы контактов}\\
 };

  \node[state,
        right = 0.5cm of mtrx1,
        align = center,
        text width = 1.5cm](img1)
  {
  \textbf{Входные \\данные}\\
};
  
  \node[state,
        right = 0.5cm of mtrx2,
        align = center,
        text width = 1.5cm](img2)
  {
  \textbf{Эталонные \\данные}
 };
 
   \node[state,
        below right = 0.2cm and 0.5cm of img1,
        align = center,
        double,
        text width = 1.5cm](nn)
  {
  \textbf{Нейронная \\сеть}
 };


  \path (grm) edge (parser)
  (sqs.north) edge (parser.west)
  (sqs) edge (pred)
  (parser) edge (mtrx1)
  (mtrx1) edge (img1)
  (pred) edge (mtrx2)
  (mtrx2) edge (img2)
  (img1.east) edge [bend left=20] (nn.north)
  (img2.east) edge [bend right=20] (nn.south)
   ;
  \end{tikzpicture}
  
 \pause
  \onslide<2>{\tikz[overlay,remember picture]{\draw[draw=red,thick,double,fill opacity=0.2, rounded corners] ($ (zzz) + (0.0,5.5)$) -- ($ (zzz) + (9.7,5.5)$) -- ($ (zzz) + (9.7,4.2)$) -- ($ (zzz) + (2.8,4.2)$) -- ($ (zzz) + (2.8,2.3)$) -- ($ (zzz) + (0.0,2.3)$) -- ($ (zzz) + (0.0,5.5)$);}}

  \onslide<3>{\tikz[overlay,remember picture]{\draw[draw=red,thick,double,fill opacity=0.2, rounded corners] ($ (zzz) + (0.0,3.5)$) rectangle ($ (zzz) + (9.7,2.2)$);}}

  \onslide<4>{\tikz[overlay,remember picture]{\draw[draw=red,thick,double,fill opacity=0.2, rounded corners] ($ (zzz) + (7.5,5.5)$) rectangle ($ (zzz) + (12.0,2.2)$);}}
\end{frame}

\begin{frame}{Нейронная сеть}

\begin{tabular}{cl}  
    \parbox{0.75\linewidth}{
        \begin{itemize}
            \item Convolutional residual neural network
                \begin{itemize}
                    \item 10 residual блоков
                    \item После них --- слой, эмулирующий алгоритм локального выравнивания последовательностей
                \end{itemize}
            \item Loss --- взвешенная попиксельная разница
            \item train/valid/test = 70\%/10\%/20\%
            %\vspace{5mm\onelineskip}
             \onslide<2> {\item Метрики
            \begin{itemize}
                \item $Precision$ --- сколько из предсказанных контактов действительно являются контактами в эталоне
                \item $Recall$ --- сколько из требуемых контактов было найдено
                \item $F1$ $score$ --- объединяющая метрика
            \end{itemize}}
        \end{itemize}
    }
    & \begin{tabular}{l}
        \includegraphics[width=2.0cm]{pics/res_unit.pdf}
    \end{tabular}  \\
\end{tabular}
    
\end{frame}

\begin{frame}{Эксперименты}
\textbf{Задачи}  
\begin{itemize}
    \item Предсказание вторичных структур цепочек тРНК с различными интервалами длин
    \item Предсказание вторичных структур цепочек тРНК с псевдоузлами
\end{itemize} 

\vspace{6mm}
\textbf{Данные}
\begin{itemize}
    \item RNAcentral (последовательности тРНК)
    \item CentroidFold (эталонные структуры)
    \item Pseudobase (цепочки и эталонные структуры с псевдоузлами)
\end{itemize}

\vspace{6mm}

\textbf{Технологии}
\begin{itemize}
    \item Платформа YaccConstructor
    \item Библиотека Keras и фреймворк Tensorflow
\end{itemize}
\vspace{6mm}

\end{frame}

\begin{frame}{Предсказание вторичных структур для цепочек с разными интервалами длин}
\begin{table}[h]
\centering
\begin{tabular}{|p{2cm}|p{2cm}||p{2cm}|p{2cm}|}
\hline
\multirow{2}{*}{Precision} & \multirow{2}{*}{85\%} & LevDist & Samples \\ \cline{3-4} 
 &  & 0 & 12\% \\ \hline
\multirow{2}{*}{Recall} & \multirow{2}{*}{89\%} & 1-5 & 19\% \\ \cline{3-4} 
 &  & 6-10 & 15\% \\ \hline
\multirow{2}{*}{FMera} & \multirow{2}{*}{86\%} & \textgreater{}10 & 16\% \\ \cline{3-4} 
 &  & error & 38\% \\ \hline
\end{tabular}
\caption{Результаты тестирования модели для длины 90}
\label{table1}
\end{table}

\end{frame}

\begin{frame}{Расширение обученных моделей на данные с псевдоузлами}
\begin{tabular}{cl}  
    \parbox{0.45\linewidth}{
Псевдоузел состоит из двух шпилек, где половина стебля одной шпильки располагается между двумя половинами стебля другой шпильки (невыразим средствами КС грамматики)
    }
    & \begin{tabular}{l}
        \includegraphics[width=5.5cm]{pics/pseudoknot.png}
    \end{tabular}  \\
\end{tabular}
\vspace{6mm}

\begin{table}[h]
\begin{tabular}{|P{2.2cm}|P{2.2cm}|P{2.2cm}|P{2.2cm}|P{2.2cm}|P{2.2cm}|}
\hline
Length & Samples & Alignment & Precision & Recall & F1 score \\ \hline \hline
50-90 & 266 & $\times$ & 74\% & 73\% & 71\% \\ \hline
\end{tabular}
\caption{Результаты тестирования моделей для данных с псевдоузлами}
\label{table2}
\end{table}

\end{frame}

\begin{frame}\frametitle{Результаты}
\begin{itemize}
    \item Разработана общая архитектура решения
    \item Проведены экспериментальные исследования на различных наборах данных
    \item Подана статья на конференцию Biata-2020
\end{itemize}
\end{frame}

\begin{frame}{Планы}
\begin{itemize}
    \item Предсказание вторичных структур для цепочек различных РНК любой длины
    \item Выбор оптимального источника эталонных данных
    \item Более тщательная разработка модели, применяющей адаптивное выравнивание.
    \item Реализация более развернутой системы тестирования результатов работы нейронных сетей
    \item Поиск новых средств, а также более тонкая настройка параметров всех моделей для улучшения результатов.
\end{itemize}

\end{frame}



\end{document}